\chapter{Riemann積分}

%\section{Riemann積分}

Riemann積分は，区間$[a,b]$上の関数$f : [a,b] \to \mathbb{R}$を積分する方法である．1854年にB.~Riemann\footnote{Riemann (1826--1866), 教授資格論文「三角級数による関数の表現可能性に関して」1854}が初めて厳密に定式化した．
この方法は，区間$[a,b]$を細かく分割して，分割した区間の上で関数$f$の値を評価し，それらの値を足し合わせることで積分を定義する．
伝統的には2つの方法がある．
一つはRiemann和による方法で，もう一つはDarboux\footnote{Darboux (1842--1917)}の上積分と下積分による方法である．
両者は同値であることが知られている．
%Riemann積分は，関数$f$が（区分的）連続である場合に特に有効である．

\section{Riemann和による定義}
\begin{definition}[区間の分割]
    区間$[a,b]$の分割$\mathcal{P}$は，区間$[a,b]$を有限個の区間へ分割したものであり，
    \[
    \mathcal{P} : a = x_0 < x_1 < \cdots < x_n = b
    \]
    と表される．分割の大きさを
    \[ |\mathcal{P}| := \max_{1 \leq i \leq n} (x_i - x_{i-1})\]
    と定義する．
\end{definition}

\begin{definition}[Riemann和]
    関数$f : [a,b] \to \mathbb{R}$の区間$[a,b]$の分割$\mathcal{P}$に対するRiemann和は，分割$\mathcal{P}$の各区間$[x_{i-1}, x_i]$において，関数$f$の値を評価する点$t_i \in [x_{i-1}, x_i]$を選び，それらの値を足し合わせたものである．
    すなわち，
    \[
        S(f,\mathcal{P},t) := \sum_{i=1}^n f(t_i)(x_i - x_{i-1}),
    \]
    ここで，$t = (t_1, t_2, \ldots, t_n)$は分割$\mathcal{P}$に対する評価点の列である．    
\end{definition}

\begin{definition}[Riemann可積分]
    関数$f : [a,b] \to \mathbb{R}$が区間$[a,b]$上でRiemann可積分であるとは，
    ある実数$I$が存在して，任意の$\varepsilon > 0$に対して，
    \red{ある$\delta > 0$が存在して，}区間$[a,b]$の\red{任意の}分割$\mathcal{P}$と評価点の列$t$に対し，
    \[
        |\mathcal P| < \delta \implies |S(f,\mathcal{P},t) - I| < \varepsilon
    \]
    が成り立つことをいう．$I$を関数$f$の区間$[a,b]$に対するRiemann積分と呼び，\[\int_a^b f(x) dx\]と表す．
\end{definition}

つまり，関数$f$のRiemann積分とは，Riemann和$S(f,\mathcal{P},t)$の細分に関する極限である．
\[
\lim_{|\mathcal{P}|\to 0} S(f,\mathcal{P},t) = \int_a^b f(x) dx.
\]
特に，関数$f$が（区分的）連続である場合は，Riemann可積分であることが知られている．
また，非有界な関数はRiemann可積分でないことも知られている（なぜか）．

\section{Darbouxの上積分と下積分による定義}

% \begin{remark}
%     関数の上限・下限を扱う際には，(a) 必要に応じて関数の有界性を仮定するか，(b) 関数の値域を\textbf{拡張実数（extended real）}に広げておく
%     \[
% \eRR :=\mathbb R\cup\{\pm\infty\}.
%     \]
%     こうすると，上限・下限が常に存在するため便利である．本章では，Darbouxの上積分・下積分の定義を述べるために限定的に拡張実数を導入する．ただし$\pm \infty$両側の拡張実数を利用する際は，$+\infty - (+\infty)$ が不定形（未定義）となるため慎重に除外する．
% \end{remark}

\begin{definition}[Darbouxの上和と下和]
    関数$f : [a,b] \to \RR$の区間$[a,b]$の分割$\mathcal{P}$に対するDarbouxの上和と下和（過剰和と不足和）は，分割$\mathcal{P}$の各区間$[x_{i-1}, x_i]$における$f$の上限と下限を用いて定義される．
    すなわち，
    \begin{align*}
        U(f,\mathcal{P}) &:= \sum_{i=1}^n (x_i - x_{i-1}) \sup_{x \in [x_{i-1}, x_i]} f(x), \\
        L(f,\mathcal{P}) &:= \sum_{i=1}^n (x_i - x_{i-1}) \inf_{x \in [x_{i-1}, x_i]} f(x).
    \end{align*}
\end{definition}

Darbouxの上和と下和は（上限と下限なので）必ずしもRiemann和と一致するわけではないが，常に
\[
L(f,\mathcal{P}) \leq S(f,\mathcal{P},t) \leq U(f,\mathcal{P})
\]
が成り立つ．

\begin{definition}[Darbouxの上積分と下積分]
    関数$f : [a,b] \to \mathbb{R}$の区間$[a,b]$の分割$\mathcal{P}$に対するDarboux上積分（下積分）は，上和（下和）の分割$\mathcal{P}$に関する下限（上限）として定義される．
    すなわち，
    \begin{align*}
        \uint{a}{b} f(x) dx &:= \inf_{\mathcal{P}} U(f,\mathcal{P}), \\
        \lint{a}{b} f(x) dx &:= \sup_{\mathcal{P}} L(f,\mathcal{P}).
    \end{align*}
\end{definition}

\begin{definition}[Darboux可積分]
    関数$f : [a,b] \to \mathbb{R}$が区間$[a,b]$上でDarboux可積分であるとは，任意の$\varepsilon > 0$に対して，区間$[a,b]$の分割$\mathcal{P}$が存在して，
    \[
        U(f,\mathcal{P}) - L(f,\mathcal{P}) < \varepsilon
    \]
    が成り立つことをいう．このとき，Darboux上積分と下積分は等しくなり，この値を関数$f$の区間$[a,b]$に対するDarboux積分と呼び，\[\dint_a^b f(x) dx\]と表す．
\end{definition}
単にDarboux上積分と下積分が等しくなることをもってDarboux可積分と定義することもある．

\begin{theorem}[Darboux]
    関数$f:[a,b] \to \mathbb{R}$が区間$[a,b]$上でRiemann可積分であることとDarboux可積分であることは同値である．
    さらに，このときRiemann積分とDarboux積分は等しい．
    \[
    \int_a^b f(x) dx = \uint{a}{b} f(x) dx = \lint{a}{b} f(x) dx = 
    \dint_a^b f(x) dx.
    \]
\end{theorem}

\section{Riemann可積分関数の例}
% スローガンとして，Riemann可積分関数 $\approxi$ 区分的に連続な有界関数である．つまり，高々有限個の点を除いて連続な関数である．
% \begin{itemize}
%     \item 区分的に連続な関数はRiemann可積分である．
%     \item 区分定数関数はRiemann可積分である．
%     \item 非有界な関数はRiemann可積分でない．
%     \item Dirichlet関数（有理数上で1，無理数上で0）はRiemann可積分でない．
%     \item Thomae関数（有理数$p/q$上で$1/q$，無理数上で0）はRiemann可積分である．
% \end{itemize}
%\paragraph{スローガン}
典型的な Riemann 可積分関数は，連続関数と区分的に連続な関数である．
ここで「区分的に連続」とは，ある分割
\[
a=x_0<x_1<\cdots<x_m=b
\]
が存在して，各開区間 $(x_{i-1},x_i)$ 上で $f$ が連続であり，
各分点 $x_i$ において有限な片側極限をもつことをいう．
後で示すように，このような関数は Riemann 可積分である．

\begin{example}
    定数関数$f(x) = c$はRiemann可積分であり，その積分値は
    \[ \int_a^b f(x) dx = c(b-a) \]
    である．
\end{example}
\begin{proof}
    区間$[a,b]$の分割$\mathcal{P}$と評価点の列$t$を任意に選ぶと，Riemann和は常に
    \[
    S(f,\mathcal{P},t) = \sum_{i=1}^n c(x_i - x_{i-1}) = c(b-a)
    \]
    となる．
    これは分割$\mathcal{P}$の選び方に依らない定数なので，$f$はRiemann可積分である．
    \[
    \int_a^b f(x)dx := \lim_{|\calP| \to 0} S(f,\calP,t) = c(b-a).
    \]
%    したがって，任意の$\varepsilon > 0$に対して，区間$[a,b]$の分割$\mathcal{P}$と評価点の列$t$が存在して，
%    \[
%    |S(f,\mathcal{P},t) - c(b-a)| = 0 < \varepsilon
%    \]
%    が成り立つ．
\end{proof}
% \begin{proof}%[Proof（Darbouxの上積分と下積分による証明）]
%     区間$[a,b]$の分割$\mathcal{P}$を任意に選ぶと，Darbouxの上和と下和は常に
%     \[
%     U(f,\mathcal{P}) = L(f,\mathcal{P}) = c(b-a)
%     \]
%     となる．したがって，任意の$\varepsilon > 0$に対して，区間$[a,b]$の分割$\mathcal{P}$が存在して，
%     \[
%     U(f,\mathcal{P}) - L(f,\mathcal{P}) = 0 < \varepsilon
%     \]
%     が成り立つ．    
% \end{proof}

% \begin{example}
%     区分定数関数は，区間$[a,b]$上の分割$\mathcal P_0 : a = x_0 < x_1 < \cdots < x_n = b$を用いて以下のように表示できる．
%     \[
%     f(x) := \sum_{i=1}^n c_i 1_{[x_{i-1},x_i)}(x), \quad f(b):= c_n.
%     \]
%     区分定数関数はRiemann可積分であり，その積分値は
%     \[ \int_a^b f(x) dx = \sum_{i=1}^n c_i (x_i - x_{i-1}) \]
%     である．
% \end{example}
% \begin{proof}
%     区間$[a,b]$の分割$\mathcal{P}$として，$\mathcal P_0$よりも細かい分割を選ぶと，Darbouxの上和と下和は常に
%     \[
%     U(f,\mathcal{P}) = L(f,\mathcal{P}) = \sum_{i=1}^n c_i (x_i - x_{i-1})
%     \]
%     となる．したがって，Darboux可積分である．よってRiemann可積分である．
    % 任意の$\varepsilon > 0$に対して，区間$[a,b]$の分割$\mathcal{P}$が存在して，
    % \[
    % U(f,\mathcal{P}) - L(f,\mathcal{P}) = 0 < \varepsilon
    % \]
    % が成り立つ．さらに，このときDarboux上積分と下積分は等しくなり，この値は
    % \[ \uint{a}{b} f(x) dx = \underline{\int_a^b} f(x) dx = \sum_{i=1}^n c_i (x_i - x_{i-1}) \]
    % である．
% \end{proof}

\begin{example}
    区間$[a,b]$上の関数$f(x) = x$はRiemann可積分であり，その積分値は
    \[ \int_a^b f(x) dx = \frac{b^2 - a^2}{2} \]
    である．
\end{example}
\begin{proof}
    
    区間$[a,b]$の分割$\mathcal{P} : a = x_0 < x_1 < \cdots < x_n = b$を用いて，
 \[
 \red{U(f,\mathcal P) - L(f,\mathcal P) = \sum_{i=1}^n (x_i-x_{i-1}
 )^2 \le (b-a)|\mathcal P| \to 0 \quad \text{as} \quad |\mathcal P| \to 0.}
 \]
 よって$f$は可積分である．一方，
    評価点の列として中点$c_i := \frac{x_i+x_{i-1}}{2}$を選ぶと，Riemann和は以下のようになる．
    \[
    S(f,\mathcal{P},c) = \sum_{i=1}^n \frac{x_i + x_{i-1}}{2} (x_i - x_{i-1}) = \frac{b^2 - a^2}{2}.
    \]
%    これは分割$\mathcal{P}$の選び方に依らない．%よってRiemann可積分である．
ここで，中点によるRiemann和$S(f,\calP,c)$とRiemann積分$\int_a^b x dx$はいずれもDarbouxの上和と下和の間にあるので，はさみうちの原理により
\[
\red{|S(f,\mathcal{P},c) - \int_a^b f(x) dx| \le U(f,\calP) - L(f,\calP) \to 0 \quad \text{as} \quad |\calP| \to 0.}
\]
    よって$\int_a^b x dx = \frac{1}{2}(b^2-a^2)$ である．
\end{proof}

\begin{example}
    区間$[a,b]$上の連続関数はRiemann可積分である．
\end{example}
\begin{proof}
    区間$[a,b]$上の連続関数$f$は一様連続である．したがって，任意の$\varepsilon > 0$と$n \in \mathbb{N}$に対して，
    ある$\delta >0$が存在して，\[|x-y| < \delta \implies |f(x) - f(y)| < \frac{\varepsilon}{b-a}\]が成り立つ．
    この$\delta$を用いて，区間$[a,b]$の分割$\mathcal{P}$として$|\mathcal P| < \delta$となるものを選ぶと，
    \[U(f,\mathcal{P}) - L(f,\mathcal{P}) = \sum_{i=1}^n \left( \sup_{x \in [x_{i-1}, x_i]} f(x) - \inf_{x \in [x_{i-1}, x_i]} f(x) \right)(x_i - x_{i-1}) \le \frac{\varepsilon}{b-a} \sum_{i=1}^n (x_i - x_{i-1}) = \varepsilon\]
    が成り立つ．ただし恒等式 $\sup_x f(x) - \inf_y f(y) = \sup_{x,y} |f(x) - f(y)|$ を使った．したがって，Riemann可積分である．
\end{proof}

% \begin{example}
%     Riemann可積分関数の和はRiemann可積分である．
% \end{example}
% \begin{proof}
%     区間$[a,b]$上のRiemann可積分関数$f$と$g$を任意に選ぶ．
%     任意の$\varepsilon > 0$に対して，区間$[a,b]$の分割$\mathcal{P}$と評価点の列$t$が存在して，
%     \[
%     |S(f,\mathcal{P},t) - \int_a^b f(x) dx| < \frac{\varepsilon}{2}, \quad |S(g,\mathcal{P},t) - \int_a^b g(x) dx| < \frac{\varepsilon}{2}
%     \]
%     が成り立つ．このとき，Riemann和の線型性により，
%     \[
%     |S(f+g,\mathcal{P},t) - \int_a^b (f(x) + g(x)) dx| < \varepsilon
%     \]
%     が成り立つ．
%     % さらに，Riemann和の積に関する性質により，
%     % \[
%     % |S(fg,\mathcal{P},t) - \int_a^b f(x)g(x) dx| < \varepsilon
%     % \]
%     % が成り立つ．
% \end{proof}


\begin{example}[区分定数関数]
区間 $[a,b]$ の分割
\[
\mathcal P_0:\ a=x_0<x_1<\cdots<x_n=b
\]
と定数 $c_1,\dots,c_n\in\mathbb R$ を用いて
\[
f(x)=\sum_{i=1}^n c_i \mathbf 1_{[x_{i-1},x_i)}(x),
\qquad
f(b)=c_n
\]
で定まる関数を区分定数関数という．
このとき $f$ は Riemann 可積分であり，
\[
\int_a^b f(x)\,dx
=
\sum_{i=1}^n c_i(x_i-x_{i-1})
\]
が成り立つ．
\end{example}

\begin{proof}
\[
J:=\sum_{i=1}^n c_i(x_i-x_{i-1}),
\qquad
M:=\max_{1\le i\le n}|c_i|
\]
とおく．任意の $\varepsilon>0$ をとる．
\[
0<\eta<\frac12\min_{1\le i\le n}(x_i-x_{i-1})
\]
かつ
\[
4M(n-1)\eta<\varepsilon
\]
となるように $\eta$ を選ぶ．
分割
\[
\mathcal P_\eta:\ 
a=x_0<x_1-\eta<x_1+\eta<\cdots<x_{n-1}-\eta<x_{n-1}+\eta<x_n=b
\]
を考える．

この分割の各小区間のうち，跳び点 $x_1,\dots,x_{n-1}$ を含まないものでは
$f$ は定数であるから，上和と下和の寄与は一致する．
一方，各跳び点 $x_i$ のまわりの区間
\[
[x_i-\eta,x_i+\eta]
\qquad
(1\le i\le n-1)
\]
では，$f$ の振れ幅は高々 $2M$，長さは $2\eta$ である．したがって
\[
U(f,\mathcal P_\eta)-L(f,\mathcal P_\eta)
\le
\sum_{i=1}^{n-1}(2\eta)\cdot 2M
=
4M(n-1)\eta
<
\varepsilon.
\]

さらに，$J$ は各跳び点近傍での寄与
\[
\eta c_i+\eta c_{i+1}
\]
をもつので，
\[
2\eta\min\{c_i,c_{i+1}\}
\le
\eta c_i+\eta c_{i+1}
\le
2\eta\max\{c_i,c_{i+1}\}
\]
より
\[
L(f,\mathcal P_\eta)\le J\le U(f,\mathcal P_\eta)
\]
が成り立つ．
したがって
\[
\underline{\int_a^b} f(x)\,dx
\le
J
\le
\overline{\int_a^b} f(x)\,dx
\]
であり，しかも
\[
\overline{\int_a^b} f(x)\,dx
-
\underline{\int_a^b} f(x)\,dx
\le
U(f,\mathcal P_\eta)-L(f,\mathcal P_\eta)
<
\varepsilon.
\]
$\varepsilon>0$ は任意であるから，上積分と下積分は一致し，その共通値は $J$ に等しい．
よって $f$ は Riemann 可積分であり，
\[
\int_a^b f(x)\,dx
=
J
=
\sum_{i=1}^n c_i(x_i-x_{i-1})
\]
となる．
\end{proof}


% \begin{example}
%     1点で不連続な\red{有界}関数$f : [a,b] \to \mathbb{R}$はRiemann可積分である．
% \end{example}
% \begin{proof}
%     $f$の不連続点を$t_1$として，区分定数関数$g : [a,b] \to \mathbb{R}$を以下のように定義する．
%     \[
% %    g(x) := f(a) 1_{[a,t_1)}(x) + f(t_1) 1_{[t_1,b]}(x).
% g(x) = \begin{cases}
%     \lim_{t \to t_1-0} f(t) & \text{if } x \in [a, t_1), \\
%     \lim_{t \to t_1+0} f(t) - \lim_{t \to t_1-0} f(t) & \text{if } x \in [t_1, b].
% \end{cases}
%     \]
%     これは区分定数関数であるからRiemann可積分である．
%     %その積分値は \[\int_a^b g(x) dx = f(a)(t_1 - a) + f(t_1)(b - t_1)\]である．さらに，
%     一方，関数$f-g$は区間$[a,b]$上の連続関数であるから，Riemann可積分である．
%     Riemann積分の線型性により，関数$f$はRiemann可積分であり，その積分値は
%     \[\int_a^b f(x) dx = \int_a^b g(x) dx + \int_a^b (f(x) - g(x)) dx\]
%     である．
% \end{proof}

% \begin{example}
%     区分的に連続な関数$f:[a,b] \to \mathbb{R}$はRiemann可積分である．
% \end{example}
% \begin{proof}
%     $f$の不連続点は有限個であるから，$1$点の場合と同様に区分的に定数となる関数$g$を定義して，$f-g$が連続関数であることを確認すればよい．
% \end{proof}
% \begin{proof}
%     $f$の不連続点を$x_1, x_2, \ldots, x_n$とする．区間$[a,b]$の分割$\mathcal{P}$として，区分的に定数となる関数$g$を定義することができる．
%     すなわち，区間$[x_{i-1}, x_i]$において，関数$f$の値を評価する点$t_i \in [x_{i-1}, x_i]$を選び，関数$g$を以下のように定義する．
%     \[
%     g(x) := \sum_{i=1}^n \left( \lim_{t \to x_i+0} f(t) - \lim_{t \to x_i-0} f(t) \right) 1_{[x_{i-1}, x_i)}(x), \quad g(b):= f(t_n).
%     \]
%     この関数$g$は区分定数関数であるからRiemann可積分であり，その積分値は
%     \[
%     \int_a^b g(x) dx = \sum_{i=1}^n f(t_i)(x_i - x_{i-1}) = S(f,\mathcal{P},t)
%     \]
%     である．さらに，関数$f-g$は区間$[a,b]$上の連続関数であるから，Riemann可積分であり，その積分値は0である．
% \end{proof}

\begin{example}[区分的に連続な関数]
ある分割
\[
a=x_0<x_1<\cdots<x_m=b
\]
が存在して，各開区間 $(x_{i-1},x_i)$ 上で $f$ が連続であり，
各分点 $x_i$ において有限な片側極限をもつとする．
このとき $f$ は Riemann 可積分である．
\end{example}

\begin{proof}
区分的連続性より，$f$ は $[a,b]$ 上で有界である．
\[
M:=\sup_{x\in[a,b]}|f(x)|<\infty
\]
とおく．任意の $\varepsilon>0$ をとる．
\[
0<\eta<\frac12\min_{1\le i\le m}(x_i-x_{i-1})
\]
かつ
\[
4Mm\eta<\frac{\varepsilon}{2}
\]
となるように $\eta$ を選ぶ．

まず
\[
[a,a+\eta],\quad
[x_1-\eta,x_1+\eta],\quad \dots,\quad
[x_{m-1}-\eta,x_{m-1}+\eta],\quad
[b-\eta,b]
\]
を「悪い部分」とする．その総延長は $2m\eta$ であるから，
この部分から来る上和と下和の差の寄与は高々
\[
2M\cdot 2m\eta=4Mm\eta<\frac{\varepsilon}{2}
\]
である．

一方，残りの閉区間
\[
[a+\eta,x_1-\eta],\quad
[x_1+\eta,x_2-\eta],\quad \dots,\quad
[x_{m-1}+\eta,b-\eta]
\]
上では $f$ は連続であるから，一様連続である．
したがって，それぞれをさらに細分して，各小区間における $f$ の振れ幅が
\[
\frac{\varepsilon}{2(b-a)}
\]
未満になるようにできる．
この細分を上の「悪い部分」と合わせた分割を $\mathcal P$ とすると，
良い部分からの寄与は
\[
\frac{\varepsilon}{2(b-a)}\cdot (b-a)=\frac{\varepsilon}{2}
\]
以下である．

以上より
\[
U(f,\mathcal P)-L(f,\mathcal P)<\varepsilon
\]
が成り立つ．
よって $f$ は Darboux 可積分であり，したがって Riemann 可積分である．
\end{proof}


% \begin{theorem}[区分求積法]
%     連続関数$f : [a,b] \to \mathbb{R}$に対して，区間$[a,b]$の$n$等分割$\mathcal{P}_n$をとり，
%     細かくしていくと，Riemann和$S(f,\mathcal{P}_n,t)$はRiemann積分$\int_a^b f(x) dx$に収束する．
%     \[
%     S(f,\mathcal{P}_n,t) = \frac{1}{n} \sum_{i=1}^n f(t_i) \to \int_a^b f(x) dx \quad (n \to \infty).
%     \]
% \end{theorem}
% \begin{proof}
%     連続関数$f$は一様連続である．したがって，任意の$\varepsilon > 0$に対して，区間$[a,b]$の分割$\mathcal{P}_n$が存在して，任意の評価点の列$t$に対して，
%     \[
%     |S(f,\mathcal{P}_n,t) - \int_a^b f(x) dx| < \varepsilon
%     \]
%     が成り立つ．
% \end{proof}

% \begin{theorem}[区分求積法]
%     Riemann可積分関数$f : [a,b] \to \mathbb{R}$に対して，区間$[a,b]$の$n$等分割$\mathcal{P}_n = \sqcup_{i=1}^n J_i^{(n)}$と，評価点の列$t^{(n)}$として$t_i^{(n)} \in J_i^{(n)}$を選ぶ．
%     % Riemann和を
%     % \[
%     % S(f,\mathcal{P}_n,t) = \frac{1}{n} \sum_{i=1}^n f(t_i)
%     % \]
%     % と定義する．
%     このとき
%     \[
%     \lim_{n \to \infty} \frac{b-a}{n} \sum_{i=1}^n f(t_i^{(n)}) = \int_a^b f(x) dx.
%     \]
% \end{theorem}
% \begin{proof}
%     区分求積によるRiemann和を$I_n := \frac{b-a}{n} \sum_{i=1}^n f(t_i^{(n)})$とおく．
%     $\mathcal P_n$によるDarbouxの上和と下和の差は
%     \[
%     U(f,\mathcal{P}_n) - L(f,\mathcal{P}_n) = \frac{1}{n} \sum_{i=1}^n \left( \sup_{x \in J_i^{(n)}} f(x) - \inf_{x \in J_i^{(n)}} f(x) \right)
%     \]
%     である．Riemann可積分関数は有界なので，この差は0に収束する．一方，
%     Riemann和は常にDarbouxの上和と下和の間にある．
%     \[
%     L(f,\mathcal{P}_n) \leq I_n \leq U(f,\mathcal{P}_n)
%     \]
%     はさみうちの原理により，Riemann和もRiemann積分に収束する．
% \end{proof}

\begin{theorem}[区分求積法]
    関数$f : [a,b] \to \mathbb{R}$はRiemann可積分とする．
    区間$[a,b]$の分割と評価点の列
    \[\mathcal{P}_k : a=x_0^{(k)} < x_1^{(k)} < \cdots < x_{n_k}^{(k)}=b, \quad x_{i_k-1}^{(k)} \le t_{i_k}^{(k)} \le x_{i_k}^{(k)}.\]
    として，$\lim_{k \to \infty} |\calP_k| = 0$ となるものを\textbf{1つ}とる
    （$|\calP_k|$ は単調減少でなくてもよい．）．
    このとき，Riemann和の列はRiemann積分に収束する．
    \[
    S(f,\calP_k,t_k) = \sum_{i=1}^{n_k} f(t_i^{(k)}) (x_{i}^{(k)} - x_{i-1}^{(k)}) \to \int_a^b f(x) dx, \quad \text{as} \quad k \to \infty.
    \]
\end{theorem}
\begin{proof}
%    各$k$毎のRiemann和を$S_k := S(f,\calP_k,t_k)$とおく．
    可積分性により，$\mathcal P_k$によるDarbouxの上和と下和の差は$0$に収束する
    \[
    U(f,\mathcal{P}_k) - L(f,\mathcal{P}_k) \to 0, \quad \text{as} \quad k \to \infty.
    \]
    % 一方，Riemann和は常にDarbouxの上和と下和の間にある．
    % \[
    % L(f,\mathcal{P}_k) \leq S(f,\calP_k,t_k) \leq U(f,\mathcal{P}_k)
    % \]
    % はさみうちの原理により，Riemann和もRiemann積分に収束する．
さて，任意の $k$ についてRiemann和は常にDarbouxの上和と下和の間にある．
\[
L(f,\mathcal P_k)\le S(f,\mathcal P_k,t^{(k)})\le U(f,\mathcal P_k)
\]
また Darboux 積分の定義から
\[
L(f,\mathcal P_k)\le I\le U(f,\mathcal P_k)
\]
も成り立つ．したがってはさみうちの原理により，
\[
\bigl|S(f,\mathcal P_k,t^{(k)})-I\bigr|
\le
U(f,\mathcal P_k)-L(f,\mathcal P_k)\to 0.
\]
ゆえに
\[
S(f,\mathcal P_k,t^{(k)})\to \int_a^b f(x)\,dx.
\]
\end{proof}

\begin{remark}
    区分求積法は，Riemann積分の定義を繰り返しているように見えるが，そうではない．Riemann可積分であることが既に分かっているときは，全ての分割$\calP$を調べる必要はなく，計算しやすいように都合よくとった１つのRiemann和列 $S(f,\calP_k,t_k)$ の収束先を調べれば，それが積分値であることを主張している．
    %証明を見れば分かる通り，分割$\calP_n$は$n$等分割でなくてもよい．また，評価点$t_i^{(n)}$は$f$の不連続点であってもよい．
\end{remark}

\begin{theorem}[微分積分学の第二基本定理]
%    関数$F:[a,b] \to \mathbb{R}$は連続かつ$(a,b)$上で微分可能であり，その微分$f := F'$は区間$[a,b]$上で連続であるとする．
$C^1$級関数$F:[a,b] \to \mathbb{R}$に対して，
    $F$の微分$f := F'$は区間$[a,b]$上でRiemann可積分であり，その積分値は
    \[
    \int_a^b f(x) dx = F(b) - F(a)
    \]
    である．
\end{theorem}
\begin{proof}
    \red{$f$は連続なので，Riemann可積分である．}
    区間$[a,b]$の分割$\mathcal{P}$を任意にとる．
    $F$に対する平均値の定理により，任意の$i$に対して，区間$[x_{i-1}, x_i]$の点$t_i$が存在して，
\[
f(t_i) = \frac{F(x_i) - F(x_{i-1})}{x_i - x_{i-1}}.
\]
この点を評価点としてRiemann和を定義すると，Riemann和$S(f,\mathcal{P},t)$は以下のようなtelescopic sumとなる．
\[
S(f,\mathcal{P},t) = \sum_{i=1}^n f(t_i)(x_i - x_{i-1}) = \sum_{i=1}^n (F(x_i) - F(x_{i-1})) = F(b) - F(a).
\]
これは分割$\calP$の取り方に依らない定数であり，特に$|\calP| \to 0$のとき$\lim_{|\calP| \to 0} S(f,\calP,t) = F(b)-F(a)$．
\red{区分求積法の議論により，}
\[
\int_a^b f(x) dx = 
\lim_{n \to 0} S(f,\mathcal{P},t) 
=F(b) - F(a).
\]
%したがって，任意の$\varepsilon > 0$に対して，区間$[a,b]$の分割$\mathcal{P}$と評価点の列$t$が存在して，
%\[
%|S(f,\mathcal{P},t) - (F(b) - F(a))| = 0 < \varepsilon
%\]が成り立つ．
\end{proof}

% \begin{theorem}
%     Riemann可積分関数$f:[a,b] \to \mathbb{R}$は有界である．
% \end{theorem}
% \begin{proof}
%     任意の$\varepsilon > 0$に対して，区間$[a,b]$の分割$\mathcal{P}$と評価点の列$t$が存在して，
%     \[
%     |S(f,\mathcal{P},t) - \int_a^b f(x) dx| < \varepsilon
%     \]
%     が成り立つ．このとき，Riemann和の定義から，任意の$i$に対して，関数$f$の値を評価する点$t_i$が存在して，
%     \[
%     |f(t_i)| < \frac{\varepsilon}{b-a} + \frac{|\int_a^b f(x) dx|}{b-a}
%     \]
%     が成り立つ．したがって，関数$f$は有界である．
% \end{proof}

\section{Riemann可積分でない例}

\begin{example}
    非有界関数はRiemann可積分でない．
\end{example}
\begin{proof}
    区間$[a,b]$上の非有界関数$f$を任意に選ぶ．非有界性の仮定より，$[a,b]$の任意の分割$\mathcal{P}$において，ある区間$[x_{i-1}, x_i]$が存在して上和または下和が非有界である．したがって，Riemann可積分でない．
\end{proof}

\begin{example}
    定義域が非有界の関数はRiemann可積分でない．
\end{example}
% \begin{proof}
%     区間$[a,\infty)$上の関数$f$を任意に選ぶ．区間$[a,b]$の分割$\mathcal{P}$を任意にとる．このとき，区間$[x_{n-1}, x_n]$に対する上和は
%     \[U(f,\mathcal{P}) \geq (x_n - x_{n-1}) \sup_{x \in [x_{n-1}, x_n]} f(x)\]
%     である．$f$は定義域が無限大であるから，ある区間$[x_{n-1}, x_n]$が存在して上和が非有界である．したがって，Riemann可積分でない．
% \end{proof}
\begin{proof}
    定義より．
\end{proof}
\begin{remark}
    Riemann積分論では，非有界空間上の関数（$f:\RR \to \RR$など）の積分として，「広義Riemann積分」という拡張概念を導入する．広義Riemann積分の一部には，Lebesgue積分よりも一般的な「積分」が含まれる．
\end{remark}

\begin{example}
    Dirichlet関数（有理数上で1，無理数上で0となる関数）$f:[0,1] \to \mathbb{R}$
    \[
    f(x) = \begin{cases}
        1 & \text{if } x \in \mathbb{Q} \cap [0,1], \\
        0 & \text{if } x \in (\mathbb{R} \setminus \mathbb{Q}) \cap [0,1].
    \end{cases}
    \]
    はRiemann可積分でない．
\end{example}
\begin{proof}
    区間$[a,b]$の分割$\mathcal{P}$と評価点の列$t$を任意にとる．
    Dirichlet関数$f$は区間$[a,b]$上で稠密に不連続点を持つ．従って，各区間$[x_{i-1}, x_i]$に対して，
    \[
    \sup_{x \in [x_{i-1}, x_i]} f(x) = 1, \quad \inf_{x \in [x_{i-1}, x_i]} f(x) = 0
    \]
    である．よってDarbouxの上和と下和は常に
    \[
    U(f,\mathcal{P}) = \sum_{i=1}^n 1 \cdot (x_i - x_{i-1}) = 1-0,
    \quad L(f,\mathcal{P}) = \sum_{i=1}^n 0 \cdot (x_i - x_{i-1}) = 0
    \]
    となる．これは分割の取り方によらないため，Darboux可積分でない．よってRiemann可積分でない．
\end{proof}

\begin{example}
    区間$[0,1]$上の有理点に番号をふる：$r : \mathbb{N} \to \mathbb{Q} \cap [0,1]$（全単射）．
    これを用いて，区間$[0,1]$上の関数列$\{f_n\}$を以下のように定義する．
    \[f_n(x) := \sum_{i=1}^n 1_{\{r(i)\}}(x).\]
    各$n$に対して，$f_n$は区分的に連続なのでRiemann可積分である．一方，極限はDirichlet関数に各点収束するから，極限関数はRiemann可積分でない．
\end{example}

\section{極限と積分の順序交換}
\begin{theorem}
    Riemann可積分関数列$\{f_n\}$が区間$[a,b]$上で一様収束して，関数$f$に収束するとする．このとき，関数$f$は区間$[a,b]$上でRiemann可積分であり，さらに，
    \[
    \lim_{n \to \infty} \int_a^b f_n(x) dx = \int_a^b f(x) dx.
    \]
\end{theorem}

\begin{proof}
    上積分と下積分の単調性により，
    \[
    \uint{a}{b} f(x) dx \le \uint{a}{b} f_n(x) dx + (b-a)\| f - f_n \|,
    \quad \lint{a}{b} f(x) dx \ge \lint{a}{b} f_n(x) dx - (b-a)\| f - f_n \|.
    \]
    $f_n$はRiemann可積分なので，%任意の$\varepsilon_1 > 0$に対して，区間$[a,b]$の分割$\mathcal{P}$と評価点の列$t$が存在して，
    \[
    \uint{a}{b} f_n(x) dx = \lint{a}{b} f_n(x) dx = \int_a^b f_n(x) dx.
    \]
    よって
    \[
    \uint{a}{b} f(x) dx - \lint{a}{b} f(x) dx \le \underbrace{\uint{a}{b} f_n(x) dx - \lint{a}{b} f_n(x) dx}_{=0} + 2(b-a)\| f - f_n \| \to 0 \quad (\text{as } n \to \infty).
    \]
    積分の線形性により，
    \[
    | \int_a^b f_n(x)dx - \int_a^b f(x)dx | \le \int_a^b dx \|f_n - f\| \to 0 \quad (\text{as } n \to \infty).
    \]
\end{proof}
% \begin{proof}
%     $\varepsilon>0$を任意にとる．
%     $f_n$が$f$に一様収束しているので，任意の$\varepsilon_1 > 0$に対して，ある$N \in \mathbb{N}$が存在して，任意の$n \geq N$に対して，任意の$x \in [a,b]$に対して，
%     \[|f_n(x) - f(x)| < \varepsilon_1\]
%     が成り立つ．特に，任意の区間$I \subset [a,b]$に対して，
%     \[
%     \sup_{t \in I} f(t) \le \sup_{t \in I}　f_n(t) + \varepsilon_1, \quad \inf_{t \in I} f(t) \ge \inf_{t \in I} f_n(t) - \varepsilon_1
%     \]
%     である．
%     各$f_n$はRiemann可積分なので，任意の$\varepsilon_2>0$に対して，区間$[a,b]$の分割$\mathcal{P}$と評価点の列$t$が存在して，
%     \[
%     U(f_n, \mathcal P) - L(f_n, \mathcal P) < \varepsilon_2, \quad |S(f_n,\mathcal{P},t) - \int_a^b f_n(x) dx| < \varepsilon_2.
%     \]
%     この$\mathcal P$と一様性を用いて，
%     \begin{align*}
%         U(f,\mathcal P) &= \sum_{j=1}^m \sup_{x \in [x_{j-1}, x_j]} f(x) (x_j - x_{j-1})\le \sum_{j=1}^m \left( \sup_{x \in [x_{j-1}, x_j]} f_n(x) + \varepsilon_1 \right) (x_j - x_{j-1}) = U(f_n, \mathcal P) + \varepsilon_1(b-a)\\
%         L(f,\mathcal P) &= \sum_{j=1}^m \inf_{x \in [x_{j-1}, x_j]} f(x) (x_j - x_{j-1})\ge \sum_{j=1}^m \left( \inf_{x \in [x_{j-1}, x_j]} f_n(x) - \varepsilon_1 \right) (x_j - x_{j-1}) = L(f_n, \mathcal P) - \varepsilon_1(b-a)
%     \end{align*}
%     したがって，
%     \[
%     U(f,\mathcal P) - L(f,\mathcal P) < U(f_n, \mathcal P) - L(f_n, \mathcal P) + 2\varepsilon_1(b-a) < \varepsilon_2 + 2\varepsilon_1(b-a).
%     \]
%     よって，関数$f$はRiemann可積分であり，その積分は以下のように評価できる．
%     \[
%     |S(f,\mathcal P,t) - \int_a^b f(x)dx| < \varepsilon_2 + 2\varepsilon_1(b-a).
%     \]
%     また，
%     \[
%     |S(f_n,\mathcal{P},t) - S(f,\mathcal{P},t)| = \sum_{j=1}^m |f_n(t_j) - f(t_j)| (x_j - x_{j-1}) < \varepsilon_1(b-a).
%     \]
%     三角不等式を用いて，$\varepsilon_2=\varepsilon/2, \varepsilon_1 = \varepsilon/(6(b-a))$とおくと，
%     \begin{align*}
%         &|\int_a^b f_n(x) dx - \int_a^b f(x)dx| \\
%         &\le |\int_a^b f_n(x) dx - S(f_n,\mathcal{P},t)| + |S(f_n,\mathcal{P},t) - S(f,\mathcal{P},t)| + |S(f,\mathcal{P},t) - \int_a^b f(x)dx| \\
%         &< \varepsilon_2 + 3\varepsilon_1(b-a) = \varepsilon.
%     \end{align*}
% \end{proof}

\begin{remark}
多くの場合に一様収束は期待できないため，より一般的な収束概念を用いて順序交換の定理を証明する必要がある．
Riemann積分の範囲では，Arzelaの優収束定理を用いて一様収束よりも弱い条件に対する順序交換の定理を証明できる．これはLebesgue積分の優収束定理の系なので，これ以上は追究しない．
\end{remark}

\section{順序交換の応用例}

\begin{theorem}[関数項級数の項別積分]
区間$[a,b]$上のRiemann可積分関数列$\{f_n\}$が定める級数$\sum_{n=1}^\infty f_n(x)$で一様収束して，関数$g$に収束するとする．このとき，関数$g$は区間$[a,b]$上でRiemann可積分であり，さらに，
    \[
    \int_a^b g(x) dx = \sum_{n=1}^\infty \int_a^b f_n(x) dx.
    \]
\end{theorem}
\begin{proof}
    関数列$\{f_n\}$の部分和を$S_N := \sum_{n=1}^N f_n$とおく．このとき，$S_N$は区間$[a,b]$上でRiemann可積分であり，さらに，
    \[
    \int_a^b S_N(x) dx = \sum_{n=1}^N \int_a^b f_n(x) dx.
    \]
    $S_N$は$g$に一様収束するから，前の定理より，
    \[
    \int_a^b g(x) dx = \lim_{N \to \infty} \int_a^b S_N(x) dx = \lim_{N \to \infty} \sum_{n=1}^N \int_a^b f_n(x) dx = \sum_{n=1}^\infty \int_a^b f_n(x) dx.
    \]
\end{proof}

例えばTaylor展開やFourier級数の項別積分などがこの定理の応用例である．

\begin{example}[一様収束しないが順序交換できる例]
    区間$[0,1]$上の関数列$\{f_n\}$を以下のように定義する．
    \[
    f_n(x) = \begin{cases}
        1 & \text{if } 0 \le x < \frac{1}{n}, \\
        0 & \text{if } \frac{1}{n} \le x \le 1.
    \end{cases}
    \]
    このとき，関数列$\{f_n\}$は区間$[0,1]$上で関数$f(x) = 0$に各点収束する．
    一方，一様収束はしない．
    ところが，任意の$n$に対して，
    \[
    \int_0^1 f_n(x) dx = \frac{1}{n} \to 0 \quad (n \to \infty), \qquad \int_0^1 f(x) dx = 0
    \]
    であるから，一様収束しないにもかかわらず，順序交換ができる．
\end{example}

\section{Riemann積分の長所と短所}

\paragraph{長所：定義が素朴であり，計算しやすい}
\begin{itemize}
    \item Riemann和 $\approx$ 区分定数関数の積分 $\approx$ 区分求積法
    \item 微分積分学の基本定理（原始関数）
\end{itemize}

\paragraph{短所：無限回操作に対して不安定}
\begin{itemize}
    \item Dirichlet関数のように不連続点が稠密に存在する関数はRiemann可積分でない．
    \item 定義域が非有界の関数はRiemann可積分でない．
    \item 非有界な関数はRiemann可積分でない．
    \item 一様収束よりも弱い収束ではRiemann可積分性が保たれない．
\end{itemize}

\paragraph{原因：Riemann積分の定義が有限和（有限分割$\mathcal P$）の極限}
なので，非有界性や無限回操作に対して不安定である．
% \[
% S(f,\mathcal P) = \sum_{i=1}^n f(t_i) (x_i - x_{i-1}), \quad \int_a^b f(x) dx = \lim_{|\mathcal P| \to 0} S(f,\mathcal P).
% \]
\[
\mathcal P : I = \bigsqcup_{i=1}^n A_i, \quad S(f,\mathcal P) = \sum_{i=1}^n f(x_i) |A_i|, \quad \int_a^b f(x) dx = \lim_{|\mathcal P| \to 0} S(f,\mathcal P).
\]

\paragraph{技術的課題：無限和（無限分割）にすると何が起きるか}
\[
\mathcal P' : I=\bigsqcup_{i=1}^\infty A_i, \quad 
S'(f,\mathcal P'):=\sum_{i=1}^\infty f(x_i) |A_i|, \quad \int_a^b f(x) dx = \lim_{|\mathcal P'| \to 0} S'(f,\mathcal P').
\]
\begin{itemize}
    \item 級数が収束するための条件は何か
    \item 分割の取り方に依らず積分が確定するための条件は何か
\end{itemize}

%なぜ極限がいるか？なぜ無限がいるか？なぜ順序交換がいるか？
% \begin{itemize}
%     \item 極限：微分積分学の基本定理を証明するためには，平均値の定理を用いてRiemann和が微分の値に収束することを示す必要がある．
%     \item 無限：Riemann積分は無限個の区間における関数値の和として定義されるため，無限が重要な役割を果たす．
%     \item 順序交換：関数列の極限と積分の順序を交換する際には，適切な収束条件が必要である．
% \end{itemize}
\paragraph{解決策：Lebesgue積分}

\begin{itemize}
    \item Riemann積分・Jordan測度を（互換性を保ちながら）修正し，より広いクラスの関数を積分可能にする
    \item 有限加法性から可算加法性へ
\end{itemize}

\paragraph{展望：何が嬉しいか}
\begin{itemize}
    \item 非有界な関数や定義域が無限大の関数も積分可能になる
    \item 不連続点が稠密に存在する関数も積分可能になる
    \item より弱い収束概念で順序交換ができるようになる
\end{itemize}

\paragraph{多くの応用先}
\begin{itemize}
    \item 確率論・数理統計学：無限回試行，確率過程
    \item 関数解析・作用素論：無限次元の「線形代数」
    \item 微分方程式論・制御理論：解空間・状態空間の無限次元性
    \item 調和解析・信号処理：Fourier級数，Fourier変換
\end{itemize}

% \paragraph{Lebesgue 積分への見通し}
% Riemann 積分では，定義域 $[a,b]$ を細かく分割し，
% 各小区間上での代表値を足し合わせることで積分を定義する．
% 言い換えると，$x$ の側を分割して，関数を区分定数関数で近似している．

% これに対して Lebesgue 積分では，関数値の側を分割し，
% \[
% A_j:=\{x\in[a,b]\mid \alpha_{j-1}\le f(x)<\alpha_j\}
% \]
% のような逆像の集合を用いて単関数
% \[
% \varphi(x)=\sum_{j=1}^m c_j \mathbf 1_{A_j}(x)
% \]
% で $f$ を近似する．
% 直感的には，Riemann 積分が「定義域を分ける積分」であるのに対し，
% Lebesgue 積分は「関数値に応じて定義域を分類する積分」である．

% この方法を厳密に定式化するためには，
% 集合 $A_j$ の「大きさ」を一般の集合に対して定める必要がある．
% そのために次章で Lebesgue 測度を導入する．

% \paragraph{何が改善されるか}
% Lebesgue 積分では，可算個の集合を自然に扱えるようになるため，
% Riemann 積分では不安定であった極限操作に対してより強い定理が成り立つ．
% その結果，非有界な関数や不連続点が多い関数も，適切な条件のもとで積分可能になる．


